\chapter{Contexto del sistema}
\label{cap:contexto}

\section{Prop\'osito y alcance}

El sistema centralizar\'a las operaciones de pr\'estamo de bienes de la Escuela. El t\'ermino \textbf{bien} comprende libros, objetos recreativos y dispositivos; evita tratar todos los recursos como equipos electr\'onicos. Cada pr\'estamo corresponde a una unidad f\'isica identificable y a un estudiante o docente responsable.

\begin{table}[htbp]
\centering
\small
\begin{tabularx}{\textwidth}{@{}>{\raggedright\arraybackslash}p{3cm}YY@{}}
\toprule
\textbf{Tipo de bien} & \textbf{Identificaci\'on y detalle} & \textbf{Consideraci\'on de uso} \\
\midrule
Libro & C\'odigo por ejemplar; t\'itulo, autor, edici\'on e ISBN si existe. & El ISBN identifica una edici\'on, no cada ejemplar. \\
Mesa de ping-pong & C\'odigo de inventario, ubicaci\'on y accesorios entregados, si corresponde. & Se propone uso en su ubicaci\'on mediante turnos con inicio y fin. \\
Visor 3D & C\'odigo de inventario; marca, modelo, serie y accesorios cuando corresponda. & Verificaci\'on del estado y componentes en entrega y devoluci\'on. \\
Parlante & C\'odigo por unidad; marca, modelo, serie y accesorios cuando corresponda. & Registro de cables u otros componentes que integren la unidad. \\
Carrito para rob\'otica & C\'odigo por unidad; modelo y componentes incluidos. & Solo unidades habilitadas por la Escuela y no adscritas a laboratorios. \\
\bottomrule
\end{tabularx}
\caption{Bienes comprendidos en el proyecto}
\end{table}

Se incluyen el registro de personas y cuentas, la consulta del cat\'alogo, el control de inventario, las reservas, la autorizaci\'on y registro de pr\'estamos, las devoluciones, las renovaciones permitidas, las garant\'ias, las incidencias, las sanciones y el historial de operaciones.

\section{Exclusiones expl\'icitas}

\begin{itemize}
    \item Pr\'estamo de computadoras de escritorio o port\'atiles.
    \item Reserva o administraci\'on de laboratorios y pr\'estamo de sus equipos, materiales o instalaciones. Esta exclusi\'on tambi\'en se aplica a un visor, parlante o carrito adscrito a un laboratorio.
    \item Otros tipos de bienes distintos de los cinco definidos, salvo una ampliaci\'on de alcance aprobada posteriormente.
    \item Pr\'estamos a personas externas y operaci\'on para otras escuelas o sedes.
    \item Procesamiento de pagos en l\'inea, integraci\'on con ERP, autenticaci\'on institucional externa y notificaciones por correo en esta versi\'on.
\end{itemize}

El estado \emph{mantenimiento} solo indicar\'a que un bien no puede prestarse; no implica desarrollar un sistema de gesti\'on de talleres o laboratorios.

\section{Usuarios, roles y perfiles}
\label{sec:actores}

Se consideran los tres roles identificados en los apuntes. Una persona puede registrarse antes de contar con acceso; para reservar o recibir un pr\'estamo debe tener una cuenta activa y un perfil habilitado.

\subsection{Administrador}

Es el responsable operativo del servicio. Registra personas y cuentas, mantiene el inventario, configura pol\'iticas aprobadas, confirma o rechaza reservas, autoriza entregas y renovaciones, registra devoluciones, administra garant\'ias e incidencias y resuelve sanciones. Puede consultar el historial general y los reportes. El rol administrativo no habilita pr\'estamos personales en esta versi\'on.

\subsection{Docente}

Puede consultar disponibilidad, solicitar reservas, recibir bienes autorizados, pedir renovaciones y consultar sus pr\'estamos, garant\'ias y sanciones. Su perfil registra c\'odigo institucional, especialidad y vinculaci\'on con la Escuela. Los l\'imites de cantidad, plazo y garant\'ia dependen de las pol\'iticas aprobadas, no de valores fijos incorporados al programa.

\subsection{Estudiante}

Puede consultar disponibilidad, solicitar reservas, recibir bienes autorizados, pedir renovaciones y consultar su historial. Su perfil registra c\'odigo universitario \'unico, programa y semestre seg\'un el plan de estudios vigente. No puede autorizar entregas, alterar el inventario ni modificar sanciones.

\section{Flujo general del servicio}

\begin{enumerate}
    \item El estudiante o docente consulta el inventario y los intervalos disponibles.
    \item Puede solicitar una reserva para una unidad y un intervalo futuro; el administrador verifica las condiciones y la confirma o rechaza.
    \item En la entrega, con o sin reserva previa, el administrador valida cuenta, perfil, disponibilidad, l\'imites, sanciones y garant\'ia exigible.
    \item El sistema registra el pr\'estamo y la pol\'itica aplicada; la unidad pasa a estar prestada. Para la mesa, la entrega representa la habilitaci\'on del uso en el lugar autorizado.
    \item Si procede una renovaci\'on, se valida el nuevo intervalo y se conserva el historial de plazos.
    \item Al finalizar el uso, el administrador registra la devoluci\'on y revisa el bien y sus accesorios. Seg\'un su condici\'on, la unidad queda disponible, en mantenimiento o de baja.
    \item Se cierra o devuelve la garant\'ia seg\'un corresponda, se registran incidencias y se determinan las sanciones aplicables. Todas las operaciones quedan en el historial.
\end{enumerate}

\section{Decisiones pendientes de validaci\'on}
\label{sec:pendientes}

Los siguientes puntos deben acordarse con la Escuela antes de implementar o poner en servicio el sistema:

\begin{itemize}
    \item Cantidad m\'axima de pr\'estamos simult\'aneos por rol, duraci\'on por tipo de bien, horarios y n\'umero de renovaciones. Los apuntes no fijan valores num\'ericos.
    \item Confirmaci\'on del uso de la mesa de ping-pong en su ubicaci\'on y de las modalidades de uso o retiro de los dem\'as bienes. Los turnos por horas son una propuesta, no una regla institucional verificada.
    \item Tipos de garant\'ia autorizados, custodia y condiciones de liberaci\'on. No se presume autorizaci\'on para retener DNI u otros documentos personales.
    \item Existencia y cuant\'ia de multas, periodos de bloqueo, tolerancias y procedimiento de apelaci\'on. No se aplicar\'an cobros sin una pol\'itica aprobada.
    \item Vigencia de perfiles, plan de estudios, responsables operativos y procedimiento ante bienes perdidos.
    \item Plataforma de despliegue, volumen de uso, respaldos y plazos de conservaci\'on de datos personales e historial.
\end{itemize}

Como simplificaci\'on de esta versi\'on se propone una sola cuenta y un solo rol por persona, una unidad por pr\'estamo y como m\'aximo una garant\'ia por pr\'estamo. El registro de varias unidades requiere operaciones individuales, todas sujetas al l\'imite total del solicitante.
